% ============================================================================
%  Capitolo 15 -- La memoria                                     (Lezione 15)
% ============================================================================
\chapter{La memoria}
\label{ch:memory}

\epigraphline{\ldots i computer e i programmi cominceranno a pensare e le
persone smetteranno.}{Dr.\ Walter Gibbs, \emph{Tron}, Steven Lisberger, 1982}

\begin{chapgoals}
Alla fine di questo capitolo saprete
\begin{itemize}
\item classificare le memorie come volatili o persistenti, in linea o fuori linea, e collocarle nella gerarchia delle memorie;
\item descrivere come SRAM e DRAM memorizzano un bit, e perché sono usate per compiti diversi;
\item spiegare che cosa sono le parole, gli indirizzi e la memoria indirizzabile al byte;
\item distinguere l'ordine little-endian da quello big-endian, e dire quale usa \riscv;
\item descrivere passo per passo come il processore legge e scrive una parola;
\item spiegare l'ingresso/uscita mappato in memoria e a che cosa serve il DMA.
\end{itemize}
\end{chapgoals}

% ----------------------------------------------------------------------------
\section{Riscaldamento sulla macchina}

\asmfile[caption={La stessa memoria letta come parola e come byte.},label=lst:mem:warm]{code/l15/s04.S}

Nel Listato~\ref{lst:mem:warm} una parola di trentadue bit occupa quattro
indirizzi consecutivi: la memoria è indirizzata byte per byte. Il byte
all'indirizzo più basso è \hex{78}, il meno significativo: l'ordine è
little-endian. \ins{lw} preleva quattro byte in una sola operazione, \ins{lbu}
ne preleva uno.

\begin{prgtry}
Eseguitelo in \prgqt e aprite la vista della memoria all'indirizzo contenuto in
\reg{t0}: i byte appaiono come \texttt{78 56 34 12}. L'ordine dei byte del
Capitolo~\ref{ch:information} è diventato un esperimento.
\end{prgtry}

Di che cosa è fatta una memoria, e che cosa succede tra il processore e la
memoria quando viene eseguita \ins{lw}?

% ----------------------------------------------------------------------------
\section{Che cos'è una memoria?}

\begin{definition}{Memoria}
La \term{memoria} è l'elemento di un calcolatore, o di un suo sottosistema,
dedicato alla persistenza dei dati e delle istruzioni dei programmi.
\end{definition}

Le memorie sono i componenti capaci di memorizzare informazione: dati,
programmi e risultati. Si classificano lungo due assi. Una memoria è
\term{volatile} se l'informazione è conservata solo finché il calcolatore è
alimentato, e \term{persistente} se è conservata anche a calcolatore spento. È
\term{in linea} se i dati sono sempre accessibili, e \term{fuori linea} se il
supporto deve essere montato prima che i dati possano essere raggiunti.

\subsection{La gerarchia delle memorie}

Nessuna tecnologia è allo stesso tempo veloce, capiente ed economica. Un
calcolatore ne usa quindi diverse, disposte in una \term{gerarchia}
(Figura~\ref{fig:mem:hier}): salendo, le memorie sono più veloci e più costose
per bit; scendendo, più capienti e meno costose per bit. I tempi di accesso
tipici vanno da meno di un nanosecondo a pochi nanosecondi per registri e
cache, decine di nanosecondi per la DRAM, decine di microsecondi per i dischi a
stato solido, millisecondi per i dischi magnetici, da secondi a minuti per il
nastro.

\begin{figure}[htb]
\centering
\begin{tikzpicture}
  \foreach \i/\n/\d/\c/\w in {0/registri/{dentro il processore}/prgred/30, 1/cache/{SRAM, sul chip del processore}/prgamber/46, 2/{memoria centrale}/DRAM/prgaccent/62, 3/{dischi a stato solido e magnetici}/{persistenti, in linea}/prgsteel/78, 4/{nastri e supporti rimovibili}/{persistenti, fuori linea}/prgviolet/94}{
    \node[blk,fill=\c,draw=\c,text=white,minimum width=\w mm,minimum height=9mm] at (0,-\i*1.05) {\textbf{\n}\\[-1pt]\scriptsize \d};
  }
  \draw[->,very thick,prgink] (5.6,-4.4) -- (5.6,0.3) node[above,lbl,align=center]{più veloci, più\\costose per bit};
  \draw[->,very thick,prgink] (-5.6,0.3) -- (-5.6,-4.4) node[below,lbl,align=center]{più capienti,\\meno costose per bit};
\end{tikzpicture}
\caption{La gerarchia delle memorie.}
\label{fig:mem:hier}
\end{figure}

% ----------------------------------------------------------------------------
\section{Un po' di storia: anelli di ferrite}

\subsection{La memoria a nuclei magnetici}

Nei primi calcolatori, prima che si diffondessero le memorie a semiconduttore,
un bit era il verso del campo magnetico di un piccolo anello di ferrite, una
ceramica magnetica (Figura~\ref{fig:mem:core}). Una memoria a nuclei è statica:
conserva il suo contenuto senza alimentazione. Le prime ricerche furono
condotte da An Wang e Way-Dong Woo a Harvard nel 1949; all'inizio degli anni
Cinquanta Jay Forrester al MIT costruì una memoria a nuclei per Whirlwind, un
calcolatore pensato come simulatore di volo.

\begin{figure}[htb]
\centering
\begin{subfigure}[b]{0.3\textwidth}\centering\includegraphics[width=\textwidth,height=38mm,keepaspectratio]{core_memory.jpg}\caption{Nuclei di ferrite e i loro fili}\label{fig:mem:core}\end{subfigure}\hfill
\begin{subfigure}[b]{0.34\textwidth}\centering\includegraphics[width=\textwidth,height=38mm,keepaspectratio]{ram_modules.jpg}\caption{Moduli di memoria di più generazioni}\label{fig:mem:modules}\end{subfigure}\hfill
\begin{subfigure}[b]{0.3\textwidth}\centering\includegraphics[width=\textwidth,height=38mm,keepaspectratio]{ram_on_board.jpg}\caption{Un modulo e il suo alloggiamento}\label{fig:mem:slot}\end{subfigure}
\caption{La memoria, ieri e oggi.}
\end{figure}

\subsection{Leggere e scrivere un nucleo}

Ogni nucleo è attraversato da fili. Due linee di indirizzo, X e Y, selezionano
un nucleo: solo il nucleo al loro incrocio riceve corrente sufficiente per
cambiare.

\emph{La lettura è distruttiva.} Il bit selezionato viene posto a 0; se
conteneva 1, compare un impulso su un terzo filo, la linea di \emph{sense}, e
altrimenti non succede niente. Il dato è stato letto e perso. La soluzione di
Wang fu che ogni lettura è seguita da una scrittura che memorizza di nuovo il
dato.

\emph{La scrittura.} Tutti i bit partono da 0. X e Y selezionano il nucleo, che
viene posto a 1; se il bit deve essere 0, una linea di \emph{inhibit} impedisce
il cambiamento. (Sense e inhibit condividono un solo filo.) Il contributo di
Forrester fu la disposizione che gestisce un gran numero di nuclei con pochi
collegamenti.

La lettura distruttiva seguita da una riscrittura è tornata oggi nella DRAM.

% ----------------------------------------------------------------------------
\section{La memoria oggi}

\subsection{SRAM e DRAM}

\begin{definition}{SRAM e DRAM}
In una \term{RAM statica} (SRAM) ogni bit è tenuto da un circuito bistabile,
come in un flip-flop. In una \term{RAM dinamica} (DRAM) ogni bit è una carica
in un minuscolo condensatore, dietro un solo transistor.
\end{definition}

Le celle di una SRAM sono disposte in una matrice e raggiunte per riga e
colonna; ciascuna richiede sei transistor (Capitolo~\ref{ch:electronics}) e
conserva il suo contenuto finché è alimentata. La SRAM è veloce e costosa, ed è
usata per le memorie cache; è anche la memoria centrale di un microcontrollore
come l'ESP32-C6, che ha 512~KB di SRAM e nessuna DRAM.

Il condensatore di una cella DRAM si scarica, e il suo contenuto deve essere
\term{rinfrescato} periodicamente da un circuito interno alla memoria. La DRAM
è economica e densa, ed è la memoria centrale di portatili, telefoni e server.
Una \emph{SDRAM} è una DRAM sincrona, il cui accesso è governato da un clock;
DDR4 e DDR5 ne sono generazioni.

\subsection{Altre memorie}

Una \emph{RAM ferroelettrica} (FeRAM) conserva il suo stato senza
alimentazione. La memoria \emph{flash} è persistente e si cancella a blocchi: è
la memoria delle chiavette USB, dei telefoni e dei dischi a stato solido,
trattati nel prossimo capitolo. Il programma di un microcontrollore vive nella
flash e le sue variabili nella SRAM.

I moduli di memoria hanno cambiato forma e dimensione negli anni
(Figura~\ref{fig:mem:modules}): SIMM, SDRAM PC100 e PC133, Rambus e DDR
contenevano da decine a centinaia di megabyte ciascuno; i moduli DDR5 di oggi
contengono da 16 a 64~GB.

\subsection{Come riconoscere la memoria}

Sulla scheda madre di un desktop o di un server la memoria sta in file di
lunghi alloggiamenti accanto al processore, che contengono i moduli
(Figura~\ref{fig:mem:slot}). Su un portatile o un telefono sta in chip saldati
vicino al processore. Sul kit \prg non c'è niente da vedere: la SRAM è dentro
l'ESP32-C6 stesso, e la flash è un piccolo chip al suo fianco.

Due parametri caratterizzano una memoria: la sua \term{dimensione}, il numero
di celle, che oggi è di gigabyte su un PC e di kilobyte su un microcontrollore;
e il suo \term{tempo di accesso}, il ritardo dopo il quale il dato è
disponibile, decine di nanosecondi per la DRAM. Per confronto, a 160~MHz un
periodo di clock dell'ESP32-C6 è di 6,25~ns.

% ----------------------------------------------------------------------------
\section{La memoria centrale}

\subsection{Parole e indirizzi}

\begin{definition}{Memoria centrale}
La \term{memoria centrale} è un array di \term{parole}; una parola è una
stringa di $m$ bit, alla quale il processore accede con una sola operazione di
lettura o di scrittura. Ogni parola è identificata da un \term{indirizzo}: un
intero compreso tra 0 e $N-1$, rappresentato su $n$ bit, con $N=2^{n}$.
\end{definition}

L'indirizzo è una parola codice, e la memoria la decodifica: è il
decodificatore del Capitolo~\ref{ch:combinational} e la piccola memoria del
Capitolo~\ref{ch:registers}, in scala maggiore. Trentadue bit di indirizzo
danno $2^{32}$ indirizzi, 4~GiB. Una memoria si dice ad \term{accesso casuale}
(RAM) quando il tempo di accesso non dipende dalla posizione del dato. I
processori accedono alla memoria con un parallelismo uguale alla lunghezza
della loro parola: 16, 32 o 64 bit.

\subsection{Memoria indirizzabile al byte}

In una memoria \term{indirizzabile al byte} il byte è la più piccola unità che
ha un indirizzo (Figura~\ref{fig:mem:bytes}). Una parola di 32 bit occupa
quattro indirizzi consecutivi, e il suo indirizzo è quello del suo primo byte.
Parole consecutive distano 4 indirizzi: è per questo che il program counter
avanza di 4, e che l'indice di un array di parole è moltiplicato per 4
(Capitolo~\ref{ch:information}).

\begin{figure}[htb]
\centering\small
\begin{tabular}{@{}r@{\quad}*{8}{|>{\centering\arraybackslash\ttfamily}p{9mm}}|@{}}
\multicolumn{1}{r}{\footnotesize indirizzo} & \multicolumn{1}{c}{\scriptsize $A$} & \multicolumn{1}{c}{\scriptsize $A{+}1$} & \multicolumn{1}{c}{\scriptsize $A{+}2$} & \multicolumn{1}{c}{\scriptsize $A{+}3$} & \multicolumn{1}{c}{\scriptsize $A{+}4$} & \multicolumn{1}{c}{\scriptsize $A{+}5$} & \multicolumn{1}{c}{\scriptsize $A{+}6$} & \multicolumn{1}{c}{}\\\cline{2-9}
 & \cellcolor{prgaccent!20}78 & \cellcolor{prgaccent!20}56 & \cellcolor{prgaccent!20}34 & \cellcolor{prgaccent!20}12 & \cellcolor{prgamber!25}48 & \cellcolor{prgamber!25}45 & \cellcolor{prgamber!25}4C & \ldots\\\cline{2-9}
\multicolumn{1}{r}{} & \multicolumn{4}{c}{\footnotesize una parola: \texttt{.word 0x12345678}} & \multicolumn{3}{c}{\footnotesize tre byte: \texttt{"HEL\ldots"}} & \multicolumn{1}{c}{}\\
\end{tabular}
\caption{Una memoria indirizzabile al byte: ogni cella è un byte.}
\label{fig:mem:bytes}
\end{figure}

\subsection{Big-endian e little-endian}

I due nomi dicono come i byte di una parola sono disposti in memoria
(Figura~\ref{fig:mem:endian}). Nell'ordine \term{little-endian} il byte meno
significativo va all'indirizzo più basso: è l'ordine di \riscv, di x86 e di
ARM. Nell'ordine \term{big-endian} il byte più significativo va all'indirizzo
più basso: è l'ordine del Motorola~68000, il processore di riferimento delle
edizioni precedenti di questo corso, e dei protocolli di Internet.

\begin{figure}[htb]
\centering\small
\begin{tabular}{@{}r@{\quad}*{4}{|>{\centering\arraybackslash\ttfamily}p{13mm}}|@{}}
\multicolumn{1}{r}{} & \multicolumn{1}{c}{\scriptsize 1000} & \multicolumn{1}{c}{\scriptsize 1001} & \multicolumn{1}{c}{\scriptsize 1002} & \multicolumn{1}{c}{\scriptsize 1003}\\\cline{2-5}
\textbf{little-endian}: \riscv, x86, ARM & \cellcolor{prggreen!25}FF & \cellcolor{prgviolet!20}E6 & \cellcolor{prgamber!25}A3 & \cellcolor{prgaccent!20}10\\\cline{2-5}
\multicolumn{5}{c}{}\\[-6pt]\cline{2-5}
\textbf{big-endian}: Motorola 68000, ordine di rete & \cellcolor{prgaccent!20}10 & \cellcolor{prgamber!25}A3 & \cellcolor{prgviolet!20}E6 & \cellcolor{prggreen!25}FF\\\cline{2-5}
\end{tabular}
\caption{La parola \hex{10A3E6FF} memorizzata all'indirizzo 1000.}
\label{fig:mem:endian}
\end{figure}

\subsection{Parole allineate e non allineate}

Una parola che comincia a un indirizzo multiplo della sua dimensione è
\term{allineata}: multipli di 4 per parole di 32 bit. Una parola allineata è
prelevata con un solo accesso alla memoria; una parola non allineata sta a
cavallo di due parole di memoria, e servono due accessi. La specifica \riscv
lascia all'implementazione il trattamento degli accessi non allineati: un load
o uno store a un indirizzo non allineato è, secondo il processore, eseguito più
lentamente o rifiutato con un'eccezione. I programmi corretti tengono allineati
i loro dati. L'assemblatore allinea i dati dichiarati con \code{.word}, e la
direttiva \code{.balign 4} forza l'allineamento dopo una stringa o un byte
(Listato~\ref{lst:mem:align}). Le istruzioni devono sempre essere allineate.

\asmfile[caption={Forzare l'allineamento di una parola.},label=lst:mem:align,numbers=none]{code/l15/s19.S}

\subsection{Parametri di una memoria}

La \term{capacità} di una memoria è il numero $N$ dei suoi registri e la
lunghezza in bit di ciascuno: 1~MB sono $2^{20}$ registri di 8 bit; «16 kW di
32 bit» sono $16\times 2^{10}$ registri di 32 bit. Il tempo di accesso è il
tempo necessario per compiere una lettura o una scrittura: il tempo per
selezionare la locazione, più il tempo per la lettura o la scrittura vera e
propria.

\begin{prgmachine}
L'ESP32-C6 ha 512~KB di SRAM per i dati, 320~KB di ROM, e una flash esterna di
alcuni megabyte per i programmi. Il menu di \prg mostra ciò che resta al
programmatore: circa 98~KB di RAM e 509~KB di flash per le cartucce. La memoria
centrale del chip è quella di otto Commodore~64, o un trentamillesimo di quella
di un portatile.
\end{prgmachine}

% ----------------------------------------------------------------------------
\section{Processore e memoria}

\subsection{L'interfaccia}

Il processore e la memoria si parlano attraverso tre bus e due registri
(Figura~\ref{fig:mem:iface}). Il \term{registro indirizzo di memoria} (MA)
contiene l'indirizzo e pilota il \term{bus indirizzi}. Il \term{registro
buffer di memoria} (MB) contiene il dato che viaggia, in un verso o
nell'altro, sul \term{bus dati}. L'unità di controllo pilota le \term{linee di
controllo}, tra le quali ci sono l'abilitazione alla lettura e l'abilitazione
alla scrittura. MA e MB sono registri interni del modello di processore del
Capitolo~\ref{ch:processor}: non sono tra i trentadue registri che un
programmatore \riscv vede.

\begin{figure}[htb]
\centering
\begin{tikzpicture}
  \node[blkacc,minimum width=44mm,minimum height=10mm] (ma) at (0,1.5) {\textbf{MA}: registro indirizzo di memoria};
  \node[blkamb,minimum width=44mm,minimum height=10mm] (mb) at (0,0) {\textbf{MB}: registro buffer di memoria};
  \node[blk,fill=prgsteel,text=white,minimum width=44mm,minimum height=9mm] (cu) at (0,-1.5) {unità di controllo};
  \node[draw=prgsteel,dashed,rounded corners,fit=(ma)(mb)(cu),inner sep=3mm,label={[lbl]above:Processore}] {};
  \node[blkvio,minimum width=26mm,minimum height=40mm] (m) at (8,0) {\textbf{Memoria}};
  \draw[->,bus,prgaccent] (ma.east) -- (ma.east-|m.west) node[midway,above,lbl,text=prgaccent]{bus indirizzi};
  \draw[<->,bus,prgamber] (mb.east) -- (mb.east-|m.west) node[midway,above,lbl,text=prgamber]{bus dati};
  \draw[->,thick,prgred] (cu.east) -- (cu.east-|m.west) node[midway,above,lbl,text=prgred]{controllo: abil.\ lettura, abil.\ scrittura};
\end{tikzpicture}
\caption{L'interfaccia tra processore e memoria.}
\label{fig:mem:iface}
\end{figure}

\subsection{Prelevare una parola: \texttt{lw t1, 0(t0)}}

\begin{table}[htb]
\centering\small
\caption{I passi di un load.}\label{tab:mem:lw}
\begin{tabularx}{\textwidth}{@{}cLl@{}}
\toprule
\textbf{Passo} & \textbf{Che cosa succede} & \textbf{Dove}\\
\midrule
1 & l'indirizzo, $\reg{t0}+0$, è calcolato dalla ALU e posto in MA & processore\\
2 & MA pilota il bus indirizzi & bus indirizzi\\
3 & l'unità di controllo attiva l'abilitazione alla lettura & linee di controllo\\
4 & la memoria decodifica l'indirizzo e mette i quattro byte della parola sul bus dati & memoria, bus dati\\
5 & il dato sul bus dati è catturato in MB & processore\\
6 & MB è copiato nel registro destinazione \reg{t1} & processore\\
\bottomrule
\end{tabularx}
\end{table}

Uno store, \code{sw t1, 0(t0)}, è simmetrico (Tabella~\ref{tab:mem:lw}):
l'indirizzo va in MA e il dato in MB, viene attivata l'abilitazione alla
scrittura, e la memoria memorizza ciò che trova sul bus dati. Il tempo di
accesso è il tempo che passa tra il passo 2 e il passo 5.

\subsection{Sincronizzare processore e memoria}

Come fa il processore a sapere che la memoria ha finito? Con la tecnica
\emph{asincrona} aspetta finché la memoria segnala, su una linea di
\term{acknowledge} (conferma), che l'operazione è completa: il processore si
adatta alla velocità della memoria, e il progetto della memoria diventa
indipendente da quello del processore, al prezzo di una maggiore complessità.
Con la tecnica \emph{sincrona} una delle linee di controllo porta gli impulsi
di un clock a frequenza fissa, e una lettura o una scrittura si completa entro
un ciclo di clock, che va quindi dimensionato sulla velocità della memoria. La
seconda tecnica è più semplice, ed è la regola dentro un singolo chip come un
microcontrollore.

% ----------------------------------------------------------------------------
\section{Memoria e ingresso/uscita}

\subsection{Collegare processore, memoria e I/O}

Il sottosistema di ingresso/uscita permette al processore di interagire con
dispositivi di ingresso (tastiera, mouse), dispositivi di uscita (schermo,
stampante) e dispositivi che sono entrambe le cose (memoria di massa, schede di
rete, porte seriali). Può essere collegato in due modi.

\begin{definition}{I/O mappato in memoria}
Con l'\term{I/O mappato in memoria} \eng{memory-mapped I/O} la memoria e i
dispositivi condividono lo stesso bus e lo stesso spazio degli indirizzi: un
registro di un dispositivo si legge e si scrive come una locazione di memoria,
e non servono istruzioni speciali.
\end{definition}

L'alternativa è l'\term{I/O isolato}: due bus fisici, o due spazi degli
indirizzi, uno per la memoria e uno per l'I/O, e istruzioni dedicate per
raggiungere i dispositivi. La famiglia Intel x86 ha entrambi: le istruzioni
\ins{in} e \ins{out} oltre ai dispositivi mappati in memoria.

\begin{prgmachine}
L'I/O mappato in memoria è la via di \riscv: \riscv non ha affatto istruzioni
di ingresso/uscita. Scrivere una parola in un registro di una periferica con
\ins{sw} è il modo in cui un programma \riscv accende un pixel o legge un
pulsante.
\end{prgmachine}

\subsection{Dove stanno le cose sull'ESP32-C6}

La Tabella~\ref{tab:mem:map} è una mappa semplificata dello spazio degli
indirizzi del chip del kit. Un solo spazio degli indirizzi serve per tutto: le
stesse trentadue linee di indirizzo raggiungono ROM, RAM, flash e ogni
dispositivo, e le stesse \ins{lw} e \ins{sw} li fanno funzionare. (La tabella è
semplificata a partire dal manuale tecnico di riferimento \citep{esp32c6trm},
che va consultato prima di fare affidamento su un indirizzo.)

\begin{table}[htb]
\centering\small
\caption{Lo spazio degli indirizzi dell'ESP32-C6, semplificato.}\label{tab:mem:map}
\begin{tabularx}{\textwidth}{@{}lLl@{}}
\toprule
\textbf{L'intervallo comincia con} & \textbf{Che cosa c'è} & \textbf{Raggiunto da}\\
\midrule
\hex{4000\ldots} & ROM: il codice che avvia il chip & prelievo delle istruzioni\\
\hex{4080\ldots} & SRAM: 512 KB per dati, stack e cartucce caricate in RAM & \ins{lw}, \ins{sw}, prelievo delle istruzioni\\
\hex{4200\ldots} & flash esterna, vista attraverso una cache: firmware e costanti & \ins{lw}, prelievo delle istruzioni\\
\hex{6000\ldots} & registri delle periferiche: GPIO, SPI, I2S, timer, Wi-Fi & \ins{lw}, \ins{sw}\\
\bottomrule
\end{tabularx}
\end{table}

L'indirizzo \hex{40800000}, al quale viene caricata ogni cartuccia
(Capitolo~\ref{ch:setup}), è l'inizio della SRAM.

\subsection{I tre bus componenti, e il DMA}

Il bus che collega processore, memoria e I/O è dunque fatto di tre bus: il
\emph{bus dati}, che trasporta i dati; il \emph{bus indirizzi}, che trasporta
l'indirizzo, di memoria o di I/O, da cui o verso cui il dato è diretto; e le
\emph{linee di controllo}: abilitazione alla lettura e alla scrittura,
acknowledge, e altre.

\begin{definition}{Accesso diretto alla memoria}
Con l'\term{accesso diretto alla memoria} (DMA, \emph{direct memory access}),
blocchi di dati passano direttamente tra un dispositivo e la memoria. Il
processore avvia il trasferimento e viene avvisato quando è terminato.
\end{definition}

Se più unità condividono il bus e possono accedere direttamente alla memoria,
un controller DMA fa da arbitro tra esse. Sull'ESP32-C6 le periferiche SPI e
I2S usano il DMA per inviare il frame buffer al display e i campioni audio
all'amplificatore: il processore non copia i pixel uno per uno.

% ----------------------------------------------------------------------------
\section{Verificate di aver capito}

\begin{quiz}{Una memoria ha 16 linee di indirizzo ed è indirizzabile al byte. Quanto è grande?}
\item 16 byte
\item 32 KiB
\item 64 KiB
\item 16 MiB
\end{quiz}

\begin{quiz}{Su \riscv, \code{value: .word 0xCAFE0001} sta all'indirizzo $A$. Che cosa legge \code{lbu t0, 0(t0)} se $\reg{t0}=A$?}
\item \hex{CA}
\item \hex{FE}
\item \hex{00}
\item \hex{01}
\end{quiz}

\begin{quiz}{Perché la DRAM va rinfrescata?}
\item Perché i suoi condensatori perdono la carica
\item Perché la lettura è lenta
\item Perché non è ad accesso casuale
\item Perché il clock è troppo veloce
\end{quiz}

\begin{quiz}{Come scrive un programma \riscv su un dispositivo di uscita?}
\item Con l'istruzione \ins{out}
\item Con \ins{sw} all'indirizzo di un registro del dispositivo
\item Con un registro speciale \reg{x32}
\item Non può: solo il firmware può farlo
\end{quiz}

\section{Esercizi}

\begin{exercise}[prima su carta, poi su \prg]
\begin{enumerate}[nosep]
\item Una memoria di 4~KiB parole di 16 bit: quante linee di indirizzo, quante
  linee di dati, quanti bit in tutto?
\item Scrivete i byte memorizzati, in ordine di indirizzo, da
  \code{.word 0xDEADBEEF} e da \code{.asciz "Ciao"} su una macchina
  little-endian.
\item Elencate i passi di \code{sw t1, 8(t0)} sull'interfaccia
  processore--memoria.
\item In \prgqt, trovate la vostra variabile contatore nella vista della
  memoria e guardate i suoi quattro byte cambiare frame per frame.
\item Memorizzate una parola con \ins{sw} e rileggetela con quattro istruzioni
  \ins{lbu}; ricostruitela con scorrimenti e \ins{or} e verificate che
  coincida.
\end{enumerate}
\expected{Su carta: 12 linee di indirizzo, 16 linee di dati, 65\,536 bit;
\texttt{EF BE AD DE} e \texttt{43 69 61 6F 00}. Su \prg: il byte meno
significativo del contatore cambia a ogni frame, il successivo ogni 256 frame.
Nell'ultimo compito i byte sono fatti scorrere di 8, 16 e 24 posizioni e poi
uniti.}
\end{exercise}

\section{Glossario}

\begin{glossarytable}
main memory & memoria centrale\\
volatile / persistent & volatile / persistente\\
magnetic core & nucleo magnetico (di ferrite)\\
refresh & rinfresco (refresh)\\
byte order & ordinamento dei byte\\
access time & tempo di accesso\\
aligned word & parola allineata\\
address / data bus & bus indirizzi / bus dati\\
read / write enable & abilitazione lettura / scrittura\\
memory-mapped I/O & I/O mappato in memoria\\
\end{glossarytable}

\section{Conclusioni e letture}

La memoria è un componente fondamentale di un calcolatore: un dispositivo
capace di conservare lo stato dei dati. Tipi diversi di memoria servono a scopi
diversi, perché nessuna tecnologia è insieme veloce, capiente ed economica. I
dati sono memorizzati come parole di bit, ciascuna con un indirizzo; \riscv
indirizza i byte e li ordina in little-endian. L'accesso avviene attraverso la
selezione (il bus indirizzi), la lettura o la scrittura (il bus dati) e le
linee di controllo. In \riscv, i dispositivi di ingresso e di uscita si
raggiungono esattamente come la memoria.

\paragraph{Letture.} \citet{patterson2020}, paragrafi 2.3 e 5.1--5.2;
\citet{stallings2021}, il capitolo sulla memoria interna; il manuale tecnico
di riferimento dell'ESP32-C6, «System and Memory» \citep{esp32c6trm}.
